\ifdefined\gkver\else\def\gkver{student}\fi
\documentclass[\gkver]{gaokaozhenti}
\begin{document}

\chapter{2019年北京理综}
%% paperType: 理综物理部分

\begin{enumerate}
\item
%% number: 13
%% paperName: 2019年北京理综第 13 题 6 分
%% typeId: 1
%% type: 单选题
%% chapter: 机械波
%% point: 波的图象
%% method: 无
%% score: 6
%% degree: 750
%% duplicateId: 0
%% body:
一列简谐横波某时刻的波形如图所示，比较介质中的三个质点 a、b、c，则\xzanswer{C}

\onepicture[w=5.5cm]{\includesvg{figs/13}}

\fourchoices[answer=C]
{此刻 a 的加速度最小}
{此刻 b 的速度最小}
{若波沿 $x$ 轴正方向传播，此刻 b 向 $y$ 轴正方向运动}
{若波沿 $x$ 轴负方向传播，a 比 c 先回到平衡位置}

%% answer: C
%% memo:
\memoanswer{
由题图可知，在简谐波的图象中，纵坐标为质点相对平衡位置的位移，在该时刻质点 a 位于正向最大位移处，由 $F=-kx$ 可知，此时质点 a 所受的回复力最大，加速度最大，故 A 错误；此时质点 b 正好处于平衡位置，加速度为零，速度最大，故 B 错误；若波沿 $x$ 轴正方向传播，由“上下坡法”可知 b 向 $y$ 轴正方向运动，故 C 正确；若波沿 $x$ 轴负方向传播，则 c 点向 $y$ 轴正方向运动，质点 a 回到平衡位置的时间等于 $\frac{1}{4}T$，而质点 c 回到平衡位置的时间小于 $\frac{1}{4}T$，故 D 错误。
}

\item
%% number: 14
%% paperName: 2019年北京理综第 14 题 6 分
%% typeId: 1
%% type: 单选题
%% chapter: 光学
%% point: 光的干涉
%% method: 无
%% score: 6
%% degree: 750
%% duplicateId: 0
%% body:
利用图 \ref{2019北京理综14a} 所示的装置（示意图），观察光的干涉、衍射现象，在光屏上得到如图 \ref{2019北京理综14b} 中甲和乙两种图样。下列关于 P 处放置的光学元件说法正确的是\xzanswer{A}

\twopicture[numstyle=arabic, labelA=2019北京理综14a, labelB=2019北京理综14b, wA=6.5cm, wB=6.5cm, align=right]
{figs/14a.png}
{figs/14b.png}

\fourchoices[answer=A]
{甲对应单缝，乙对应双缝}
{甲对应双缝，乙对应单缝}
{都是单缝，甲对应的缝宽较大}
{都是双缝，甲对应的双缝间距较大}

%% answer: A
%% memo:
\memoanswer{
在单色光的双缝干涉实验中，干涉条纹等间距分布；而单缝衍射实验中，衍射条纹中间宽，向两侧逐渐变窄，所以乙图中 P 处为双缝，甲图中 P 处为单缝，故 A 正确，BCD 错误。
}

\item
%% number: 15
%% paperName: 2019年北京理综第 15 题 6 分
%% typeId: 1
%% type: 单选题
%% chapter: 气体性质
%% point: 物体的内能
%% method: 无
%% score: 6
%% degree: 750
%% duplicateId: 0
%% body:
下列说法正确的是\xzanswer{A}

\fourchoices[answer=A]
{温度标志着物体内大量分子热运动的剧烈程度}
{内能是物体中所有分子热运动所具有的动能的总和}
{气体压强仅与气体分子的平均动能有关}
{气体膨胀对外做功且温度降低，分子的平均动能可能不变}

%% answer: A
%% memo:
\memoanswer{
温度反映了物体内分子做无规则运动的剧烈程度，故 A 正确；内能是物体中所有分子热运动所具有的动能和分子势能的总和，故 B 错误；微观上看，气体压强既与气体分子平均动能有关，也与气体分子的密度有关，故 C 错误；温度是气体分子平均动能的量度，温度降低，平均动能减小，故 D 错误。
}

\item
%% number: 16
%% paperName: 2019年北京理综第 16 题 0 分
%% typeId: 1
%% type: 单选题
%% chapter: 磁场
%% point: 带电粒子在磁场中的运动
%% method: 无
%% score: 0
%% degree: 550
%% duplicateId: 0
%% body:
如图所示，正方形区域内存在垂直纸面的匀强磁场。一带电粒子垂直磁场边界从 a 点射入，从 b 点射出。下列说法正确的是\xzanswer{C}

\onepicture[w=4.2cm]{\includesvg{figs/16}}

\fourchoices[answer=C]
{粒子带正电}
{粒子在 b 点速率大于在 a 点速率}
{若仅减小磁感应强度，则粒子可能从 b 点右侧射出}
{若仅减小入射速率，则粒子在磁场中运动时间变短}

%% answer: C
%% memo:
\memoanswer{
由左手定则可知，带正电粒子所受洛伦兹力向上，只有粒子带负电其运动方向才能向下，故 A 错误；由于洛伦兹力不做功，a 点速率与 b 点速率相等，故 B 错误；粒子运动的半径 $R=\frac{mv}{qB}$，仅减小磁感应强度 $B$，$R$ 增大，则粒子有可能从 b 点右侧射出，故 C 正确；若 $v$ 减小，则 $R$ 减小，由图可知，粒子转过的圆心角变大，又由粒子在磁场中运动的周期 $T=\frac{2\pi m}{qB}$，$t=\frac{\theta}{2\pi}\cdot T$ 知，$\theta$ 增大，粒子在磁场中运动时间变长，故 D 错误。

\onepicture[w=3.5cm]{figs/16_an.jpg}
}

\item
%% number: 17
%% paperName: 2019年北京理综第 17 题 6 分
%% typeId: 1
%% type: 单选题
%% chapter: 电场
%% point: 电场线
%% method: 无
%% score: 6
%% degree: 650
%% duplicateId: 0
%% body:
如图所示，a、b 两点位于以负点电荷 $-Q$（$Q>0$）为球心的球面上，c 点在球面外，则\xzanswer{D}

\onepicture[w=4.5cm]{\includesvg{figs/17}}

\fourchoices[answer=D]
{a 点场强的大小比 b 点大}
{b 点场强的大小比 c 点小}
{a 点电势比 b 点高}
{b 点电势比 c 点低}

%% answer: D
%% memo:
\memoanswer{
设该球面半径为 $R$，a 点和 b 点处于同一球面上，半径相等，则由 $E=k\frac{Q}{r^2}$ 可知，a 点和 b 点的场强大小相同，c 点在球面外，则 c 点场强小于 a、b 点场强，故 AB 错误；a 点和 b 点在同一等势面上，电势相等，故 C 错误；由沿电场线方向电势降低可知，c 点电势高于 a 点和 b 点，故 D 正确。
}

\item
%% number: 18
%% paperName: 2019年北京理综第 18 题 6 分
%% typeId: 1
%% type: 单选题
%% chapter: 曲线运动
%% point: 人造地球卫星
%% method: 无
%% score: 6
%% degree: 650
%% duplicateId: 0
%% body:
2019 年 5 月 17 日，我国成功发射第 45 颗北斗导航卫星，该卫星属于地球静止轨道卫星（同步卫星）。该卫星\xzanswer{D}

\fourchoices[answer=D]
{入轨后可以位于北京正上方}
{入轨后的速度大于第一宇宙速度}
{发射速度大于第二宇宙速度}
{若发射到近地圆轨道所需能量较少}

%% answer: D
%% memo:
\memoanswer{
同步卫星只能定点于赤道上方固定高度处，故 A 错误；绕地球做圆周运动的卫星，由万有引力提供向心力得，$G\frac{Mm}{r^2}=m\frac{v^2}{r}$，得 $v=\sqrt{\frac{GM}{r}}$，则 $r$ 越大，$v$ 越小，第一宇宙速度是近地卫星运行的最大速度，即同步卫星速度小于第一宇宙速度，故 B 错误；同步卫星发射速度要小于第二宇宙速度大于第一宇宙速度，故 C 错误；卫星轨道越低，发射时所需要达到的速度越小，所需的能量越少，故 D 正确。
}

\item
%% number: 19
%% paperName: 2019年北京理综第 19 题 0 分
%% typeId: 1
%% type: 单选题
%% chapter: 光学
%% point: 光电效应
%% method: 无
%% score: 0
%% degree: 650
%% duplicateId: 0
%% body:
光电管是一种利用光照射产生电流的装置，当入射光照在管中金属板上时，可能形成光电流。表中给出了 6 次实验的结果。

\begin{center}
\small
\setlength{\tabcolsep}{4pt}
\begin{tabular}{|c|c|c|c|c|c|}
\hline
组 & 次 & 入射光子的能量/$\UeV$ & 相对光强 & 光电流大小/$\UmA$ & 逸出光电子的最大动能/$\UeV$ \\
\hline
第一组 & 1 & 4.0 & 弱 & 29 & 0.9 \\
\hline
 & 2 & 4.0 & 中 & 43 & 0.9 \\
\hline
 & 3 & 4.0 & 强 & 60 & 0.9 \\
\hline
第二组 & 4 & 6.0 & 弱 & 27 & 2.9 \\
\hline
 & 5 & 6.0 & 中 & 40 & 2.9 \\
\hline
 & 6 & 6.0 & 强 & 55 & 2.9 \\
\hline
\end{tabular}
\end{center}

由表中数据得出的论断中不正确的是\xzanswer{B}

\fourchoices[answer=B]
{两组实验采用了不同频率的入射光}
{两组实验所用的金属板材质不同}
{若入射光子的能量为 $5.0\UeV$，逸出光电子的最大动能为 $1.9\UeV$}
{若入射光子的能量为 $5.0\UeV$，相对光强越强，光电流越大}

%% answer: B
%% memo:
\memoanswer{
两组实验中，入射光子的能量不同，所以两组实验采用了不同频率的入射光，故 A 错误；

由光电效应方程 $E_k=h\nu-W$ 可知，两组实验所用金属板的逸出功均为 $3.1\UeV$，材质相同，故 B 正确；

由光电效应方程 $E_k=h\nu-W$ 可得入射光电子能量为 $5.0\UeV$ 时，逸出光电子的最大初动能为 $5.0\UeV-3.1\UeV=1.9\UeV$，故 C 错误；

分析两组实验可知，入射光子频率不变，入射光越强，光电流越大，故 D 错误。
}

\item
%% number: 20
%% paperName: 2019年北京理综第 20 题 6 分
%% typeId: 1
%% type: 单选题
%% chapter: 运动和力
%% point: 力学单位制
%% method: 无
%% score: 6
%% degree: 450
%% duplicateId: 0
%% body:
国际单位制（缩写 SI）定义了米（m）、秒（s）等 7 个基本单位，其他单位均可由物理关系导出。例如，由 m 和 s 可以导出速度单位 $\Umdsn[a]$。历史上，曾用“米原器”定义米，用平均太阳日定义秒。但是，以实物或其运动来定义基本单位会受到环境和测量方式等因素的影响，而采用物理常量来定义则可避免这种困扰。1967 年用铯-133 原子基态的两个超精细能级间跃迁辐射的频率 $\Delta\nu=9\,192\,631\,770\UHz$ 定义 s；1983 年用真空中的光速 $c=299\,792\,458\Umdsn$ 定义 m。2018 年第 26 届国际计量大会决定，7 个基本单位全部用基本物理常量来定义（对应关系如图，例如，s 对应 $\Delta\nu$，m 对应 $c$）。新 SI 自 2019 年 5 月 20 日（国际计量日）正式实施，这将对科学和技术发展产生深远影响。下列选项不正确的是\xzanswer{D}

\onepicture[w=4.5cm]{figs/20.png}

\fourchoices[answer=D]
{7 个基本单位全部用物理常量定义，保证了基本单位的稳定性}
{用真空中的光速 $c$（$\Umdsn[a]$）定义 m，因为长度 $l$ 与速度 $v$ 存在 $l=vt$，而 s 已定义}
{用基本电荷 $e$（$\UC[a]$）定义安培（$\UA[a]$），因为电荷量与电流 $I$ 存在 $I=q/t$，而 s 已定义}
{因为普朗克常量 $h$（$\UJcdots[a]$）的单位中没有 $\Ukg[a]$，所以无法用它来定义质量单位}

%% answer: D
%% memo:
\memoanswer{
由题意可得，7 个基本单位都用物理常量定义，故 A 不符合题意；选项 BC 正是长度和电流的定义方式，故 BC 不符合题意；$1\UJcdots[a]=1\UNcdotm[a]\cdot\Us[a]=1\Ukg[a]\Um^2\Usn$，则普朗克常量 $h$（$\UJcdots[a]$）的单位中含有 $\Ukg[a]$，可以定义质量，故 D 正确。
}

\item
%% number: 21
%% paperName: 2019年北京理综第 21 题 18 分
%% typeId: 4
%% type: 实验题
%% chapter: 曲线运动
%% point: 研究平抛运动实验
%% method: 无
%% score: 18
%% degree: 450
%% duplicateId: 0
%% body:
用如图 \ref{2019北京理综21a} 所示装置研究平抛运动。将白纸和复写纸对齐重叠并固定在竖直的硬板上。钢球沿斜槽轨道 PQ 滑下后从 Q 点飞出，落在水平挡板 MN 上。由于挡板靠近硬板一侧较低，钢球落在挡板上时，钢球侧面会在白纸上挤压出一个痕迹点。移动挡板，重新释放钢球，如此重复，白纸上将留下一系列痕迹点。

\begin{enumerate}
	\item
	下列实验条件必须满足的有\tkanswer[2cm]{BD}。

	A．斜槽轨道光滑

	B．斜槽轨道末段水平

	C．挡板高度等间距变化

	D．每次从斜槽上相同的位置无初速度释放钢球

	\item
	为定量研究，建立以水平方向为 $x$ 轴、竖直方向为 $y$ 轴的坐标系。

	a．取平抛运动的起始点为坐标原点，将钢球静置于 Q 点，钢球的\tkanswer[2cm]{球心}（选填“最上端”、“最下端”或者“球心”）对应白纸上的位置即为原点；在确定 $y$ 轴时\tkanswer[1.5cm]{需要}（选填“需要”或者“不需要”）$y$ 轴与重锤线平行。

	b．若遗漏记录平抛轨迹的起始点，也可按下述方法处理数据：如图 \ref{2019北京理综21b} 所示，在轨迹上取 A、B、C 三点，AB 和 BC 的水平间距相等且均为 $x$，测得 AB 和 BC 的竖直间距分别是 $y_1$ 和 $y_2$，则 $\frac{y_1}{y_2}$\tkanswer[1cm]{大于}$\frac{1}{3}$（选填“大于”、“等于”或者“小于”）。可求得钢球平抛的初速度大小为\tkanswer[2.6cm]{$x\sqrt{\frac{g}{y_2-y_1}}$}（已知当地重力加速度为 $g$，结果用上述字母表示）。

	\item
	为了得到平抛物体的运动轨迹，同学们还提出了以下三种方案，其中可行的是\tkanswer[1.2cm]{AB}。

	A．从细管水平喷出稳定的细水柱，拍摄照片，即可得到平抛运动轨迹

	B．用频闪照相在同一底片上记录平抛小球在不同时刻的位置，平滑连接各位置，即可得到平抛运动轨迹

	C．将铅笔垂直于竖直的白纸板放置，笔尖紧靠白纸板，铅笔以一定初速度水平抛出，将会在白纸上留下笔尖的平抛运动轨迹

	\item
	伽利略曾研究过平抛运动，他推断：从同一炮台水平发射的炮弹，如果不受空气阻力，不论它们能射多远，在空中飞行的时间都一样。这实际上揭示了平抛物体\tkanswer[1.5cm]{B}。

	A．在水平方向上做匀速直线运动

	B．在竖直方向上做自由落体运动

	C．在下落过程中机械能守恒

	\item
	牛顿设想，把物体从高山上水平抛出，速度一次比一次大，落地点就一次比一次远，如果速度足够大，物体就不再落回地面，它将绕地球运动，成为人造地球卫星。

	同样是受地球引力，随着抛出速度增大，物体会从做平抛运动逐渐变为做圆周运动，请分析原因。

	\tkanswer[12cm]{物体初速度较小时，运动范围很小，引力可以看作恒力——重力，做平抛运动；随着物体初速度增大，运动范围变大，引力不能再看作恒力；当物体初速度达到第一宇宙速度时，做圆周运动而成为地球卫星。}
\end{enumerate}

\twopicture[numstyle=arabic, labelA=2019北京理综21a, labelB=2019北京理综21b, align=right, wA=8cm, wB=6.5cm]
{figs/21a.png}
{figs/21b.png}

%% answer: 
\jdanswer{
\begin{enumerate}
	\item
	BD

	\item
	球心，需要，大于，$v_0=x\sqrt{\frac{g}{y_2-y_1}}$

	\item
	AB

	\item
	B

	\item
	物体初速度较小时，运动范围很小，引力可以看作恒力——重力，做平抛运动；随着物体初速度增大，运动范围变大，引力不能再看作恒力；当物体初速度达到第一宇宙速度时，做圆周运动而成为地球卫星。

\end{enumerate}
}
%% memo:
\memoanswer{
（1）实验中应保证每次小球滑出轨道后要平抛，所以斜槽轨道末段必须水平，选项 B 正确；同时要求每次抛出时的水平速度相同，故需要每次释放小球的高度相同，但不必苛求轨道光滑，故 A 错误、D 正确；挡板是为了确定平抛运动的轨迹，不必等间距变化，故 C 错误。

（2）a．小球不能看作质点，初始位置应是球心在竖直平板上水平的投影点，实验中 $y$ 轴记录小球竖直方向上的位移，必须与重垂线平行。

b．若 A 为起始点，则 $\frac{y_1}{y_2}=\frac{\frac{1}{2}gT^2}{\frac{1}{2}g(2T)^2-\frac{1}{2}gT^2}=\frac{1}{3}$。现在 A 点不是起始点，过 A 点时 $v_y\neq0$，由 $y=\overline{v_y}t$ 可知比值变大。由题图可知竖直方向 $y_2-y_1=g\Delta t^2$，水平方向有 $x=v_0\Delta t$，解得 $v_0=x\sqrt{\frac{g}{y_2-y_1}}$。

（3）水柱足够细，水面下降缓慢，水流出时的初速度几乎相同，故 A 可行；B 选项用频闪照片的方式获取数据更准确，故 B 可行；铅笔抛出时所受空气阻力较大，同时伴有转动，故 C 不可行。

（4）由题意可知，炮弹竖直方向上的高度相同，如果不受空气阻力，发射后落地时间相同，揭示的是平抛物体竖直方向上的运动是自由落体运动，故 B 正确。

（5）速度较小时，重力的方向可认为不变，依据平抛的规律，速度增大，落点变远；但地球表面不是平面，当速度很大时，需要考虑到地球的形状近似为球形，物体运动轨迹逐渐趋近圆，当物体受到的引力刚好等于向心力时，物体将绕地心做圆周运动，而不返回地面。
}

\item
%% number: 22
%% paperName: 2019年北京理综第 22 题 16 分
%% typeId: 6
%% type: 计算题
%% chapter: 电磁感应
%% point: 电磁感应电路分析
%% method: 无
%% score: 16
%% degree: 650
%% duplicateId: 0
%% body:
如图所示，垂直于纸面的匀强磁场磁感应强度为 $B$。纸面内有一正方形均匀金属线框 abcd，其边长为 $L$，总电阻为 $R$，ad 边与磁场边界平行。从 ad 边刚进入磁场直至 bc 边刚要进入的过程中，线框在向左的拉力作用下以速度 $v$ 匀速运动，求：

\begin{enumerate}
	\item
	感应电动势的大小 $E$；
	\item
	拉力做功的功率 $P$；
	\item
	ab 边产生的焦耳热 $Q$。
\end{enumerate}

\onepicture[w=6cm, align=right]{\includesvg{figs/22}}

%% answer: 
\jdanswer{
\begin{enumerate}
	\item
	$E=BLv$

	\item
	$P=\frac{B^2L^2v^2}{R}$

	\item
	$Q=\frac{B^2L^3v}{4R}$

\end{enumerate}
}
%% memo:
\memoanswer{
（1）由法拉第电磁感应定律可得，感应电动势为 $E=BLv$。

（2）线框中的感应电流为 $I=\frac{E}{R}$，线框做匀速直线运动，拉力大小等于安培力大小 $F=BIL$，故拉力的功率为 $P=Fv=\frac{B^2L^2v^2}{R}$。

（3）线框 ab 边电阻为 $R_{ab}=\frac{R}{4}$，此过程所用时间为 $t=\frac{L}{v}$，ab 边产生的焦耳热为 $Q=I^2R_{ab}t=\frac{B^2L^3v}{4R}$。
}

\item
%% number: 23
%% paperName: 2019年北京理综第 23 题 18 分
%% typeId: 6
%% type: 计算题
%% chapter: 电路
%% point: 电容
%% method: 无
%% score: 18
%% degree: 450
%% duplicateId: 0
%% body:
电容器作为储能器件，在生产生活中有广泛的应用。对给定电容值为 $C$ 的电容器充电，无论采用何种充电方式，其两极间的电势差 $u$ 随电荷量 $q$ 的变化图像都相同。

（1）请在图 \ref{2019北京理综23a} 中画出上述 $u-q$ 图像。类比直线运动中由 $v-t$ 图像求位移的方法，求两极间电压为 $U$ 时电容器所储存的电能 $E_p$。

（2）在如图 \ref{2019北京理综23b} 所示的充电电路中，$R$ 表示电阻，$E$ 表示电源（忽略内阻）。通过改变电路中元件的参数对同一电容器进行两次充电，对应的 $q-t$ 曲线如图 \ref{2019北京理综23c} 中①②所示。

a．①②两条曲线不同是\tkanswer[1cm]{$R$}（选填 $E$ 或 $R$）的改变造成的；

b．电容器有时需要快速充电，有时需要均匀充电。依据 a 中的结论，说明实现这两种充电方式的途径：

\tkanswer[12cm]{减小电阻 $R$，可以实现对电容器更快速充电；增大电阻 $R$，可以实现更均匀充电。}

（3）设想使用理想的“恒流源”替换（2）中电源对电容器充电，可实现电容器电荷量随时间均匀增加。请思考使用“恒流源”和（2）中电源对电容器的充电过程，填写下表（选填“增大”、“减小”或“不变”）。

\begin{center}
\begin{tabular}{|c|c|c|}
\hline
 & “恒流源” & （2）中电源 \\
\hline
电源两端电压 & & \\
\hline
通过电源的电流 & & \\
\hline
\end{tabular}
\end{center}

\threepicture[numstyle=arabic, labelA=2019北京理综23a, labelB=2019北京理综23b, labelC=2019北京理综23c, align=right, wA=4cm, wB=4.5cm, wC=6cm]
{\includesvg{figs/23a}}
{figs/23b.png}
{figs/23c.png}

%% answer: 
\jdanswer{
\begin{enumerate}
	\item
	见解析

	\item
	$R$，减小电阻 $R$，可以实现对电容器更快速充电；增大电阻 $R$，可以实现更均匀充电。

	\item
	电源两端电压：使用“恒流源”时增大，使用（2）中电源时不变；通过电源的电流：使用“恒流源”时不变，使用（2）中电源时减小。

\end{enumerate}
}
%% memo:
\memoanswer{
（1）$u-q$ 图线如图所示。

\onepicture[w=4cm]{figs/23_an.png}

电压为 $U$ 时，电容器带电 $Q$，图线和横轴围成的面积为所储存的电能 $E_p$，则 $E_p=\frac{1}{2}QU$，又 $Q=CU$，故 $E_p=\frac{1}{2}CU^2$。

（2）a．$R$。

b．减小电阻 $R$，可以实现对电容器更快速充电；增大电阻 $R$，可以实现更均匀充电。

（3）填表如下：使用“恒流源”充电时，电源两端电压增大，通过电源的电流不变；使用（2）中电源充电时，电源两端电压不变，通过电源的电流减小。
}

\item
%% number: 24
%% paperName: 2019年北京理综第 24 题 20 分
%% typeId: 6
%% type: 计算题
%% chapter: 运动和力
%% point: 变加速直线运动
%% method: 无
%% score: 20
%% degree: 350
%% duplicateId: 0
%% body:
雨滴落到地面的速度通常仅为几米每秒，这与雨滴下落过程中受到空气阻力有关。雨滴间无相互作用且雨滴质量不变，重力加速度为 $g$。

\begin{enumerate}
	\item
	质量为 $m$ 的雨滴由静止开始，下落高度 $h$ 时速度为 $u$，求这一过程中克服空气阻力所做的功 $W$。

	\item
	将雨滴看作半径为 $r$ 的球体，设其竖直落向地面的过程中所受空气阻力 $f=kr^2v^2$，其中 $v$ 是雨滴的速度，$k$ 是比例系数。

	a．设雨滴的密度为 $\rho$，推导雨滴下落趋近的最大速度 $v_m$ 与半径 $r$ 的关系式；

	b．示意图中画出了半径为 $r_1$、$r_2$（$r_1>r_2$）的雨滴在空气中无初速下落的 $v-t$ 图线，其中\tkanswer[1cm]{①}对应半径为 $r_1$ 的雨滴（选填①、②）；若不计空气阻力，请在图中画出雨滴无初速下落的 $v-t$ 图线。

	\item
	由于大量气体分子在各方向运动的几率相等，其对静止雨滴的作用力为零。将雨滴简化为垂直于运动方向面积为 $S$ 的圆盘，证明：圆盘以速度 $v$ 下落时受到的空气阻力 $f\propto v^2$（提示：设单位体积内空气分子数为 $n$，空气分子质量为 $m_0$）。
\end{enumerate}

\onepicture[w=6cm, align=right]{\includesvg{figs/24}}

%% answer: 
\jdanswer{
\begin{enumerate}
	\item
	$W=mgh-\frac{1}{2}mu^2$

	\item
	a．$v_m=\sqrt{\frac{4\pi r\rho g}{3k}}$；b．①，图见解析。

	\item
	$f\propto nm_0Sv^2$

\end{enumerate}
}
%% memo:
\memoanswer{
（1）雨滴由静止开始下落，由动能定理得 $mgh-W=\frac{1}{2}mu^2$，可得这一过程克服空气阻力做功为 $W=mgh-\frac{1}{2}mu^2$。

（2）a．雨滴下落的过程，由牛顿第二定律得 $mg-f=ma$，解得 $a=g-\frac{kr^2v^2}{m}$，当加速度为零时，雨滴趋近于最大速度 $v_m$，又雨滴质量 $m=\frac{4}{3}\pi r^3\rho$，由 $a=0$，可得雨滴最大速度 $v_m=\sqrt{\frac{4\pi r\rho g}{3k}}$。

b．①，如图所示。

\onepicture[w=5.5cm]{figs/24_an.png}

（3）根据题设条件：大量气体分子在各方向运动的几率相等，其对静止雨滴的作用力为零。以下只考虑雨滴下落的定向运动。简化的圆盘模型如图。

\onepicture[w=5cm]{figs/24_an2.png}

设空气分子与圆盘碰撞前后相对速度大小不变。在 $\Delta t$ 时间内，与圆盘碰撞的空气分子质量为 $\Delta m=Sv\Delta tnm_0$，以 $F$ 表示圆盘对气体分子的作用力，由动量定理得 $F\Delta t\propto\Delta m\times v$，解得 $F\propto nm_0Sv^2$，由牛顿第三定律可知圆盘所受空气阻力 $f\propto v^2$，采用不同的碰撞模型，也可得到相同结论。
}

\end{enumerate}

\end{document}
